%! Semi-log Plots — Unit 5, Lesson 5.8 — Lesson Plan
\documentclass[10pt]{article}
\usepackage{precalculus-article}
\usepackage{precalculus-boxes}
\usepackage{graphicx}
\graphicspath{{images/}}
\newcommand{\TallMath}[1]{$\displaystyle #1\rule[-1.4em]{0pt}{3.2em}$}
\newcommand{\CourseName}{Precalculus}
\newcommand{\SchoolYear}{2026--2027}
\newcommand{\MeetingLength}{55 minutes}
\newcommand{\UnitNumberName}{Unit 5: Logarithmic Functions \quad}
\newcommand{\LessonNumberName}{Lesson 5.8: Semi-log Plots}

\begin{document}
\begin{center}
    {\Large\bfseries \CourseName: \SchoolYear \\
    {\small\bfseries \UnitNumberName \LessonNumberName}}
\end{center}

\begin{tcolorbox}[colback=sky, colframe=navy, boxrule=0.9pt, arc=2mm,
    left=2mm, right=2mm, top=1.5mm, bottom=1.5mm]
    \textbf{Primary Objective:} Students determine whether exponential models are appropriate
    by examining semi-log plots (where exponential data appears linear), apply linear function
    techniques to semi-log graphs, and construct the linearization of exponential data using
    the formula $\log_n(y) = (\log_n b)\cdot x + \log_n a$.
    \hfill\textit{\small Skills: AP 2.B, AP 3.C \quad LO 2.15.A, LO 2.15.B}
\end{tcolorbox}

\renewcommand{\arraystretch}{1.6}
\begin{skillbox}[Priority Ideas \& Skills]{goldbox}
\begin{tabularx}{\linewidth}{@{}X X@{}}
\textbf{AP Skills} &
\textbf{Key Understandings} \\
\midrule
\textbf{AP 2.B} --- Construct equivalent representations. &
\begin{itemize}[leftmargin=1.2em,itemsep=2pt,topsep=0pt]
  \item A semi-log plot has one logarithmically scaled axis; exponential data appears linear
        on such a plot. (EK 2.15.A.1)
  \item For $y = ab^x$, the linearized form is
        $\log_n y = (\log_n b)\cdot x + \log_n a$. (EK 2.15.B.2)
\end{itemize} \\
\textbf{AP 3.C} --- Support conclusions with evidence. &
\begin{itemize}[leftmargin=1.2em,itemsep=2pt,topsep=0pt]
  \item Linearity on a semi-log plot supports an exponential model. (EK 2.15.A.1)
  \item No constant needs to be added to $y$-values to reveal exponential structure.
        (EK 2.15.A.2)
  \item Slope $= \log_n b$ and intercept $= \log_n a$ recover the model. (EK 2.15.B.1)
\end{itemize} \\
\end{tabularx}
\end{skillbox}

\begin{skillbox}[{Vocabulary, Concepts, \& Theorems}]{greenbox}
\renewcommand{\arraystretch}{1.4}
\begin{tabularx}{\linewidth}{@{} >{\leavevmode\bfseries\color{navy}}p{0.33\linewidth} X @{}}
Semi-log plot & A graph in which one axis (typically the $y$-axis) is logarithmically scaled
                so that equal spacing represents equal ratios (powers of the base) rather than
                equal differences. \\
Logarithmic scale & An axis scaled so tick marks appear at $n^0, n^1, n^2, \ldots$; equal
                   spacing represents equal \emph{multiplicative} change. \\
Linearization & The transformation $\log_n(y) = (\log_n b)\cdot x + \log_n a$ that converts
                the exponential model $y = ab^x$ into a linear function of $x$. \\
Slope of linearized data & In the linearized model, slope $= \log_n b$; recover the base via
                           $b = n^{\text{slope}}$. \\
$y$-intercept of linearized data & In the linearized model, intercept $= \log_n a$; recover
                                   the initial value via $a = n^{\text{intercept}}$. \\
\end{tabularx}
\end{skillbox}

\begin{skillbox}[Activate Prior Knowledge \& Spiral Review]{sky}
    Spiral review (warm-up) covers: computing $y$-values for $y = 3\cdot2^x$ at
    $x = 0,1,2,3,4$; taking $\log_{10}$ of those values and observing a constant additive
    increment of $\approx 0.301$; identifying that constant increment as the slope of
    $\log_{10}(y)$ vs.\ $x$.  These connect directly to EK 2.15.B.1--2 and prime students
    for the linearization derivation.
\end{skillbox}

\begin{skillbox}[Hook]{sky}
    Pose to the class: \textbf{``A scientist plots bacterial population data on a regular
    graph and on a semi-log graph.  On the regular graph the curve bends sharply upward.
    On the semi-log graph the same data forms a perfectly straight line.  What does this
    tell the scientist?''}

    Let pairs discuss 1 min, then cold-call.  Guide toward: a straight line on a semi-log
    graph means the data grows exponentially, because applying $\log$ converts multiplicative
    growth into additive growth, which plots as a line.  This motivates the entire lesson.
\end{skillbox}

\begin{skillbox}[Lesson]{sky}
\begin{multicols}{2}
\textbf{Part 1 (10 min): What Is a Semi-log Plot? (EK 2.15.A.1--2)}

Describe a semi-log plot: the $y$-axis is scaled so equal spacing represents powers of 10
($10^0, 10^1, 10^2, \ldots$).  Exponential data falls on a straight line; non-exponential
data curves.

Advantage (EK 2.15.A.2): unlike the Lesson 2.6 approach of adding a constant to $y$-values
before checking for proportional differences, a semi-log plot requires no constant addition.
Linearity of $\log_{10}(y)$ vs.\ $x$ is sufficient evidence.

\columnbreak

\textbf{Part 2 (15 min): Linearizing $y = ab^x$ (EK 2.15.B.2)}

Take $\log_{10}$ of both sides:
\[
  \log_{10}(y) = \log_{10}(a) + x\log_{10}(b).
\]
So $\log_{10}(y) = (\log_{10} b)\cdot x + \log_{10} a$ is linear in $x$ with slope
$m = \log_{10} b$ and intercept $c = \log_{10} a$.

Worked example: $y = 5\cdot3^x$:
\[
  \log_{10}(y) \approx 0.477\,x + 0.699.
\]

\textbf{Part 3 (10 min): Recovering the Exponential Model (EK 2.15.B.1)}

From slope $m$ and intercept $c$: $b = 10^m$ and $a = 10^c$.

Example: slope $= 0.602$, intercept $= 0.699$:
$b = 10^{0.602}\approx 4$, $a = 10^{0.699}\approx 5$; model: $y = 5\cdot4^x$.

\textbf{Part 4 (5 min): Connecting to Lesson 2.6}

Semi-log plots bypass the need to add a constant to $y$-values (EK 2.15.A.2).  If
$\log_{10}(y)$ vs.\ $x$ is linear, the model is exponential --- no constant required.
\end{multicols}
\end{skillbox}

\begin{skillbox}[Active Monitoring]{redbox}
\begin{itemize}[leftmargin=1.2em,itemsep=2pt,topsep=0pt]
  \item After Part 1: ``What would \emph{linear} data (not exponential) look like on a
        semi-log plot?'' --- expect a curve, not a line.
  \item During Part 2: watch for $\log_{10}(a\cdot b^x) = (\log_{10}a)(\log_{10}b^x)$
        (product-of-logs error); emphasize the log-of-product rule.
  \item After Part 2: ``If slope of $\log_{10}(y)$ vs.\ $x$ is negative, what does that
        say about $b$?'' --- $b < 1$; model is exponential decay.
  \item During activity: verify students compute $b = 10^m$ and not $b = m$ or $b = \log m$.
\end{itemize}
\end{skillbox}

\begin{skillbox}[Group Work \& Differentiation]{redbox}
\begin{itemize}[leftmargin=1.2em,itemsep=2pt,topsep=0pt]
  \item \textbf{Tier R (Remediate):} Given drug-concentration data $C = 100,50,25,12.5,6.25$
        at $t = 0,1,2,3,4$, complete the $\log_{10}(C)$ table, observe constant rate of
        change $\approx -0.301$, and conclude an exponential model is appropriate.
  \item \textbf{Tier A (Apply):} Write the linearized model $\log_{10}(C) = mt + c$; recover
        $a$ and base $b$ via $10^m$ and $10^c$; write and verify $C(t) = 100\cdot(0.5)^t$;
        predict $C(6)$.
  \item \textbf{Tier E (Extension):} Apply linearization to an imperfect data set where the
        rate of change of $\log_{10}(C)$ is not constant; compute a best-fit slope; recover
        an approximate exponential model; compare data quality between both sets.
\end{itemize}
\end{skillbox}

\begin{skillbox}[Individual Work \& Assessment]{redbox}
Exit ticket (2 questions, 1 page):
\begin{enumerate}[leftmargin=1.5em,itemsep=2pt,topsep=2pt]
  \item Table: $x = 0,1,2,3$, $y = 4, 12, 36, 108$.  Compute $\log_{10}(y)$ for each $x$.
        Find slope and $y$-intercept of $\log_{10}(y)$ vs.\ $x$.
        (Tests EK 2.15.B.2 --- identifying $\log_{10} b$ and $\log_{10} a$.)
  \item From slope and intercept, identify $a$ and $b$ and write $y = a\cdot b^x$.
        Verify with $y(1)$.
        (Tests EK 2.15.B.1 --- recovering the exponential model.)
\end{enumerate}

Sort by result: students missing \#1 need practice with $\log_{10}$ computation and the
linearization formula; students missing \#2 need the inverse step $b = 10^m$.
\end{skillbox}

\begin{skillbox}[Reinforcement \& Extension]{goldbox}
\textbf{Homework (6 problems):} conceptual explanation of why exponential data linearizes on
a semi-log plot (EK 2.15.A.1); linearizing $y = 2\cdot5^x$; recovering an exponential model
from given slope and intercept; four-point data table with full linearization and model
recovery; advantage of semi-log over adding constants (EK 2.15.A.2); AP-style MC identifying
the slope of a linearized model.  Extension: repeat with natural log (change-of-base
comparison).

\textbf{Preview of Lesson 2.16:} Applying regression to linearized data --- using technology
to find the line of best fit for $(x,\,\log_{10}(y))$ pairs, then converting back to an
exponential model.
\end{skillbox}

\end{document}
